\chapter*{Abstract}
\addcontentsline{toc}{chapter}{Abstract}

Thin film solar cells require efficient light-trapping layers to maximise broadband absorption, but using poorly scaling full-field
electromagnetic solvers bottlenecks the development of novel nanophotonic advances. Traditional gradient-based topology optimization
can easily be trapped in local minima. To overcome this challenge, this thesis explores the uses of machine learning in the design
process of 2D patterned grayscale light-trapping layers across multiple high refractive materials (Si, TiO$_2$, Si$_3$N$_4$.) A robust data generation pipeline
in python was developed and validated using Rigorous Coupled Wave Analysis (RCWA). Three distinct neural network surrogate models
were trained to map light-trapping geometries to optical absorptance spectra, bypassing the expensive field calculations entirely. Compared to the classical RCWA solver,
these networks are multiple orders of magnitude faster, while remaining precise, at an average error of only about $2.5\%$. The networks are able to upscale wavelength resoltion at virtually no extra time complexity.
Surrogate based optimization was tested and showed great promise in replacing or complementing classical methods. This work demonstrates that deep learning may provide a robust alternative to broadband inverse design problems.